\section{Conclusion}\label{sec:conc}
We defined FGFV, a core calculus for an object-oriented language
with symmetric multiple dispatch, multiple inheritance, and
parametric polymorphism with variance.
In order to ensure type soundness of FGFV, it must guarantee no ambiguous
method invocations at run time, and preserve type even when a method
invocation is dispatched to a different method at run time from the method
chosen at compile time.
Thus, we formalized overloading rules using existential types and
universal types quantified over method type parameters to mechanically
check whether overloaded polymorphic methods are valid.
As a result, in a well-formed program that satisfies the overloading rules,
there always exists the single most specific method to invoke
with respect to both static argument types and dynamic argument types.
Further, with the static types of method invocation expressions
delivered to run time, the dynamic semantics of FGFV that uses the new
dynamic dispatch algorithm we presented is proven to make
type-sound choices among applicable method declarations.
